Дана сторона квадрата $a$. Найти квадрат его периметра $P = (4 \cdot a)^2$.
Дана сторона квадрата $a$. Найти его площадь $S = a^2$.
Даны стороны прямоугольника $a$ и $b$. Найти его площадь $S = a \cdot b$ и периметр $P = 2 \cdot (a + b)$.
Дан диаметр окружности $d$. Найти ее длину $L = \pi \cdot d$. В качестве значения $\pi$ использовать $3.14$.
Дана длина ребра куба $a$. Найти объем куба $V = a^3$ и площадь его поверхности $S = 6 \cdot a^2$.
Даны длины ребер $a, b, c$ прямоугольного параллелепипеда. Найти его объем $V = a \cdot b \cdot c$ и площадь поверхности $S = 2 \cdot (a \cdot b + b \cdot c + a \cdot c)$.
Найти длину окружности $L$ и площадь круга $S$ заданного радиуса $R$: $L = 2 \cdot \pi \cdot R$, $S = \pi \cdot R^2$. В качестве значения $\pi$ использовать $3.14$.
Даны два числа $a$ и $b$. Найти их среднее арифметическое: $(a + b)/2$.
Даны два неотрицательных числа $a$ и $b$. Найти их среднее геометрическое, то есть квадратный корень из их произведения: $\sqrt{a \cdot b}$.
Даны два ненулевых числа. Найти сумму, разность, произведение и частное их квадратов.
Даны два ненулевых числа. Найти сумму, разность, произведение и частное их модулей.
Даны катеты прямоугольного треугольника $a$ и $b$. Найти его гипотенузу $c = \sqrt{a^2 + b^2}$ и периметр $P = a + b + c$.
Даны два круга с общим центром и радиусами $R_1$ и $R_2$ ($R_1 > R_2$). Найти площади этих кругов $S_1$ и $S_2$, а также площадь $S_3$ кольца, внешний радиус которого равен $R_1$, а внутренний радиус равен $R_2$: $S_1 = \pi \cdot (R_1)^2$, $S_2 = \pi \cdot (R_2)^2$, $S_3 = S_1 - S_2$. В качестве значения $\pi$ использовать $3.14$.
Дана длина $L$ окружности. Найти ее радиус $R$ и площадь $S$ круга, ограниченного этой окружностью, учитывая, что $L = 2 \cdot \pi \cdot R$, $S = \pi \cdot R^2$. В качестве значения $\pi$ использовать $3.14$.
Дана площадь $S$ круга. Найти его диаметр $D$ и длину $L$ окружности, ограничивающей этот круг, учитывая, что $L = 2 \cdot \pi \cdot R$, $S = \pi \cdot R^2$. В качестве значения $\pi$ использовать $3.14$.
Найти расстояние между двумя точками с заданными координатами $x_1$ и $x_2$ на числовой оси: $|x_2 - x_1|$.
Даны три точки $A$, $B$, $C$ на числовой оси. Точка $C$ расположена между точками $A$ и $B$. Найти длины отрезков $AC$ и $BC$ и их сумму.
Даны три точки $A$, $B$ и $C$ на числовой оси. Точка $C$ расположена между точками $A$ и $B$. Найти произведение длин отрезков $AC$ и $BC$.
Даны координаты двух противоположных вершин прямоугольника: $(x_1, y_1)$, $(x_2, y_2)$. Стороны прямоугольника параллельны осям координат. Найти периметр и площадь данного прямоугольника.
Поменять местами содержимое переменных $A$ и $B$ и вывести новые значения $A$ и $B$.
Даны переменные $A$, $B$, $C$. Изменить их значения, переместив содержимое $A$ в $B$, $B$ --- в $C$, $C$ --- в $A$, и вывести новые значения переменных $A$, $B$, $C$.
Найти значение функции $y = 3x^6 - 6x^2 - 7$ при данном значении $x$.
Найти значение функции $y = 4(x - 3)^6 - 7(x - 3)^3 + 2$ при данном значении $x$.
Дано значение угла $\alpha$ в градусах ($0 < \alpha < 360$). Определить значение этого же угла в радианах, учитывая, что $180^\circ = \pi$ радианов. В качестве значения $\pi$ использовать $3.14$.
Дано значение угла $\alpha$ в радианах ($0 < \alpha < 2 \cdot \pi$). Определить значение этого же угла в градусах, учитывая, что $180^\circ = \pi$ радианов. В качестве значения $\pi$ использовать $3.14$.
Найти значение функции $y = 6x^2 - 3x - 7$ при данном значении $x$.
Подсчитайте, сколько очков набрала команда «Динамо» в первом круге чемпионата России по хоккею, если известно, что $m$ встреч она выиграла, $n$ встреч проиграла, $k$ встреч закончились ничьими, полагая, что за выигрыш команда получает 2 очка, за ничью -- 1 очко, за проигрыш -- 0 очков.
Розничная цена костюма составляет $R$ рублей. Торговое наложение магазина составляет $T\%$ от оптовой цены. Определите оптовую цену костюма.
Оптовая цена костюма составляет $P$ рублей. Торговое наложение магазина составляет $T\%$ от оптовой цены. Определите розничную цену костюма.
Найти значение функции $y = 5(x - 2)^5 - x + 4$ при данном значении $x$.
